\section{Related Work}
\label{sec:related-work}

Related work falls along a spectrum of filesystem ownership. Namespace
mechanisms construct views before lookup, and access-control hooks mediate or
observe operations near the VFS. \namei puts bounded object-selection policy in
VFS name resolution, while FUSE, stackable filesystems, custom filesystems, and
metadata services require broader filesystem-service or object-model ownership.

\subsection{FUSE Policy Services}

FUSE~\cite{linux_fuse}, FUSE passthrough~\cite{linux_fuse_passthrough},
ExtFUSE~\cite{extfuse}, and FUSE-BPF~\cite{fuse_bpf_patch} demonstrate the
value of placing filesystem policy in, or accelerating, a filesystem-service
boundary. They are the direct context for RQ2 because they can implement the
same path-view policy under the same source oracle, but keep the policy inside a
FUSE or stackable-filesystem request path.

\namei instead keeps policy needed by the source oracle at VFS name resolution
when the workload only needs path selection or visibility over existing
objects.

\subsection{Custom, Stackable, And Metadata-Service Filesystems}

Wrapfs~\cite{wrapfs}, Bento~\cite{bento}, and custom filesystems demonstrate
broader filesystem-extension boundaries. They can implement richer behavior and,
in Bento's case, reduce the risk of in-kernel filesystem development. They are
the direct context for RQ3 because the developer still owns filesystem methods,
storage semantics, or framework-specific filesystem behavior. These mechanisms
remain the right boundary when the workload needs custom file operations,
synthetic contents, persistence, or richer semantics.

DeltaFS, IndexFS, and TableFS show that full metadata services and specialized
filesystems are valuable when the problem is metadata distribution, indexing,
or storage layout~\cite{deltafs_sc21,indexfs,tablefs}. They are not the main
workload baselines for \namei, but they clarify the boundary because such
workloads require a metadata service or full filesystem when the oracle needs
distributed metadata or a new filesystem object model.

\subsection{Namespace Construction}

Bind mounts, projected volumes, symlink forests, copy trees, and OverlayFS are
practical ways to present alternate views~\cite{kubernetes_projected_volumes,linux_overlayfs}.
They preserve ordinary kernel filesystem behavior, but express policy by
constructing or updating namespace state. \namei asks when the policy can stay
at lookup and directory-enumeration time instead of materializing a view before
the transition completes.

\subsection{Access Control And Observation Hooks}

Landlock, BPF LSM programs, and fanotify show that Linux already exposes
specialized hooks around filesystem access and events~\cite{linux_landlock,linux_bpf_lsm,linux_fanotify}.
Among neighboring mechanisms, these hooks sit between materialized namespace
construction and \namei. They can run verified policy in the kernel, but their
natural role is to deny, mediate, or observe operations. \namei selects which
existing object a name denotes during resolution and then returns control to the
ordinary VFS and lower filesystem. Access-control denial remains outside the
current action set because the primary claim is name-resolution selection and
visibility.

\subsection{Source Workload Systems}

Agent filesystems are closest to our workloads, but they are workload sources
rather than direct cost baselines~\cite{agentfs_repo,branchfs_repo,sandlock_repo,yolofs_repo}.
They show that agent workspaces need branch, sandbox, staging, and
snapshot-like views, while their implementations also own behavior beyond
\namei's boundary.
Here, custom or stackable filesystem ownership provides the RQ3 boundary
context, while FUSE remains the RQ2 cost comparison. \namei applies when the
same source oracle can be satisfied by changing path visibility or object
selection while leaving the rest of that ownership with the source system or
lower filesystem.

Environment and cache systems provide a second workload source, not a filesystem
mechanism baseline. SWE-Factory-Gym, MEnvData-SWE, and SWE-rebench V2 provide
build/test oracles for environment state and cache reuse
policies~\cite{swe_factory_repo,menvagent_repo,swe_rebench_v2_repo}. Agent
runtimes and service/config systems similarly motivate workload state, but their
runtime orchestration, environment synthesis, projected-volume materialization,
and reload behavior remain outside \namei's mechanism claim.
